\chapter{Brouwer-Type Fixed-Point Ideas}
\label{ch:ii17}

Brouwer-type fixed-point ideas replace metric contraction by topology. Existence can survive without uniqueness or an iterative rate, revealing a distinct mechanism from Banach's theorem.

\section*{Learning goals}
The reader should be able to use the central definitions, prove the structural results, diagnose the role of each hypothesis, and solve representative analytic problems.

\section*{Conceptual roadmap}
\[
\boxed{\text{definition}\;\longrightarrow\;\text{estimate}\;\longrightarrow\;\text{theorem}\;\longrightarrow\;\text{application}.}
\]

\section{Fixed points without contraction}
Continuous maps can have fixed points for topological rather than metric reasons.

\section{Interval case}
The intermediate-value theorem already yields a one-dimensional Brouwer theorem.

\section{Convex compact sets}
Brouwer's theorem applies to continuous self-maps of compact convex subsets of Euclidean space.

\section{No-retraction intuition}
A ball cannot continuously retract onto its boundary.

\section{Degree intuition}
Topological degree counts preimages with orientation and obstructs disappearance of fixed points.

\section{Existence versus uniqueness}
Brouwer gives existence but generally no uniqueness.

\section{Applications}
Equilibria in games, economics, and nonlinear systems often use Brouwer-type existence.

\section{Comparison with Banach}
Banach is quantitative and metric; Brouwer is qualitative and topological.

\section{Core structural results}

\begin{theorem}[Interval fixed-point theorem]\label{thm:ii17-01}
Every continuous $f:[a,b]\to[a,b]$ has a fixed point.
\begin{proof}
Apply the intermediate-value theorem to $g(x)=f(x)-x$: $g(a)\ge0$ and $g(b)\le0$.
\end{proof}
\end{theorem}

\begin{theorem}[Brouwer fixed-point theorem]\label{thm:ii17-02}
Every continuous map from a nonempty compact convex subset of $\mathbb R^n$ to itself has a fixed point.
\begin{proof}
The full proof uses a topological obstruction such as no-retraction, Sperner's lemma, or degree; the key point is that a fixed-point-free map would induce impossible boundary data.
\end{proof}
\end{theorem}

\begin{theorem}[No-retraction consequence]\label{thm:ii17-03}
There is no continuous retraction from the closed ball onto its boundary sphere.
\begin{proof}
A retraction would contradict Brouwer by composing radial constructions to obtain a fixed-point-free self-map of the ball.
\end{proof}
\end{theorem}

\section{Worked examples}

\begin{example}[Interval proof]\label{ex:ii17-worked-01}
Prove the fixed-point theorem on $[0,1]$. For $g=f-\mathrm{id}$, $g(0)\ge0$ and $g(1)\le0$; continuity gives a zero.
\end{example}

\begin{example}[Failure without compactness]\label{ex:ii17-worked-02}
Give a continuous self-map of $(0,1)$ with no fixed point. $f(x)=(x+1)/2$ maps $(0,1)$ to itself and has only fixed point $1$, outside the domain.
\end{example}

\begin{example}[Failure without self-map condition]\label{ex:ii17-worked-03}
Why does $f(x)=x+1$ on $[0,1]$ not contradict Brouwer? Its image is not contained in $[0,1]$.
\end{example}

\begin{example}[Nonuniqueness]\label{ex:ii17-worked-04}
Give a Brouwer map with many fixed points. The identity map fixes every point.
\end{example}

\section{Solved dossiers}
The dossiers are longer solved problems. Legacy mappings guide topic selection where explicit retained source blocks exist; otherwise fresh canonical problems complete the chapter.

\begin{problem}[Interval proof]\label{prob:ii17-dossier-01}
Prove the fixed-point theorem on $[0,1]$.
\end{problem}
\begin{solution}
For $g=f-\mathrm{id}$, $g(0)\ge0$ and $g(1)\le0$; continuity gives a zero.
\end{solution}

\begin{problem}[Failure without compactness]\label{prob:ii17-dossier-02}
Give a continuous self-map of $(0,1)$ with no fixed point.
\end{problem}
\begin{solution}
$f(x)=(x+1)/2$ maps $(0,1)$ to itself and has only fixed point $1$, outside the domain.
\end{solution}

\begin{problem}[Failure without self-map condition]\label{prob:ii17-dossier-03}
Why does $f(x)=x+1$ on $[0,1]$ not contradict Brouwer?
\end{problem}
\begin{solution}
Its image is not contained in $[0,1]$.
\end{solution}

\begin{problem}[Nonuniqueness]\label{prob:ii17-dossier-04}
Give a Brouwer map with many fixed points.
\end{problem}
\begin{solution}
The identity map fixes every point.
\end{solution}

\begin{problem}[Banach versus Brouwer]\label{prob:ii17-dossier-05}
Which theorem gives an iteration rate?
\end{problem}
\begin{solution}
Banach does; Brouwer generally supplies existence only.
\end{solution}

\begin{problem}[Convexity role]\label{prob:ii17-dossier-06}
Why is an annulus not covered directly by Brouwer's convex-set statement?
\end{problem}
\begin{solution}
An annulus is not convex; straight segments between points may leave it.
\end{solution}

\begin{problem}[A rotation of a disk]\label{prob:ii17-dossier-07}
Why must a disk rotation have a fixed point?
\end{problem}
\begin{solution}
The center is fixed, consistent with Brouwer.
\end{solution}

\begin{problem}[Sperner intuition]\label{prob:ii17-dossier-08}
How does Sperner's lemma relate to Brouwer?
\end{problem}
\begin{solution}
A labeled triangulation forces a fully labeled simplex; shrinking mesh sizes produce approximate fixed points whose compactness limit is an exact one.
\end{solution}

\begin{problem}[Equilibrium interpretation]\label{prob:ii17-dossier-09}
How can a fixed point encode equilibrium?
\end{problem}
\begin{solution}
A response map sends a state to its updated best-response or normalized state; a fixed point is unchanged by the update and thus represents equilibrium.
\end{solution}

\begin{problem}[Topological obstruction]\label{prob:ii17-dossier-10}
What does 'no retraction' mean intuitively?
\end{problem}
\begin{solution}
One cannot continuously push every interior point to the boundary while leaving boundary points fixed.
\end{solution}

\begin{problem}[Approximate fixed points]\label{prob:ii17-dossier-11}
Why are compactness arguments natural in Brouwer proofs?
\end{problem}
\begin{solution}
Combinatorial or discretized arguments produce approximate fixed points; compactness extracts a convergent subsequence and continuity turns the limiting approximation error into zero.
\end{solution}

\begin{problem}[Brouwer synthesis]\label{prob:ii17-dossier-12}
What fundamentally distinguishes Brouwer from Banach?
\end{problem}
\begin{solution}
Brouwer relies on topology of compact convex sets and gives existence; Banach relies on quantitative shrinking in complete metric spaces and gives uniqueness plus convergence.
\end{solution}

\section{Exercises with complete solutions}

\begin{exercise}\label{exr:ii17-01}
Does every continuous $[0,1]\to[0,1]$ map have a fixed point?
\end{exercise}
\begin{hint}
Yes.
\end{hint}
\begin{solution}
Yes.
\end{solution}

\begin{exercise}\label{exr:ii17-02}
Can it have several?
\end{exercise}
\begin{hint}
Identity.
\end{hint}
\begin{solution}
Yes.
\end{solution}

\begin{exercise}\label{exr:ii17-03}
Does Brouwer guarantee uniqueness?
\end{exercise}
\begin{hint}
No.
\end{hint}
\begin{solution}
No.
\end{solution}

\begin{exercise}\label{exr:ii17-04}
Does compactness matter?
\end{exercise}
\begin{hint}
Open interval counterexample.
\end{hint}
\begin{solution}
Yes.
\end{solution}

\begin{exercise}\label{exr:ii17-05}
Does convexity appear in the standard theorem?
\end{exercise}
\begin{hint}
Yes.
\end{hint}
\begin{solution}
Yes.
\end{solution}

\begin{exercise}\label{exr:ii17-06}
Is a disk convex?
\end{exercise}
\begin{hint}
Straight segments stay inside.
\end{hint}
\begin{solution}
Yes.
\end{solution}

\begin{exercise}\label{exr:ii17-07}
Is a circle convex?
\end{exercise}
\begin{hint}
Chord leaves circle.
\end{hint}
\begin{solution}
No.
\end{solution}

\begin{exercise}\label{exr:ii17-08}
What one-dimensional theorem proves the interval case?
\end{exercise}
\begin{hint}
Intermediate value theorem.
\end{hint}
\begin{solution}
IVT.
\end{solution}

\section*{Chapter summary}
The chapter's tools now form part of the canonical Volume II analysis toolkit and will be used in later approximation, fixed-point, and convergence arguments.
